% !TEX root = ../thesis.tex
% ==============================================================================
%  CHAPTER 6: CONCLUSIONS & FUTURE WORK (فصل ششم: نتیجه‌گیری و کارهای آینده) (chapters/chapter6_conclusions.tex)
% ==============================================================================
%  This chapter summarizes the primary findings and contributions of the
%  thesis, addresses theoretical and computational limitations, and outlines
%  promising directions for future research.
% ==============================================================================

\chapter{نتیجه‌گیری و کارهای آتی}\label{chap:chapter6_conclusions}

این فصل به جمع‌بندی جامع مسیر طی‌شده در این پایان‌نامه، از صورت‌بندی صوری مسئله‌ی ممیزی مدل‌های زبانی بزرگ تا طراحی، پیاده‌سازی و ارزیابی تجربی چارچوب پیشنهادی اختصاص دارد. در ادامه، ابتدا دستاوردهای روش‌شناختی و تجربی پژوهش در پرتو نتایج حاصل از ارزیابی‌های تطبیقی در برابر خطوط پایه‌ی جعبه‌سفید و جعبه‌سیاه بازخوانی و صورت‌بندی می‌شوند. سپس، محدودیت‌های ساختاری، معرفت‌شناختی و محاسباتی حاکم بر مسئله‌ی انتساب در سیستم‌های مسدود واکاوی شده و در نهایت، افق‌های تحقیقاتی نوین برای گسترش و کاربست این رهیافت در بهینه‌سازی سامانه‌ها و ممیزی‌های کاربردی ترسیم می‌گردند.


\section{جمع‌بندی دستاوردها و نتایج}\label{sec:6-1}

هدف بنیادین این پژوهش، توسعه‌ی یک متدولوژی پساعملی و مستقل از مدل برای انتساب ویژگی در مدل‌های زبانی بزرگ با رژیم دسترسی جعبه‌سیاه بود؛ سناریویی که در آن تحلیل‌گر صرفاً به متن محقق‌شده‌ی ورودی و خروجی دسترسی دارد و پارامترها، گرادیان‌ها و حتی توزیع احتمالات خروجی (لوجیت‌ها) در پسِ رابط‌های تجاری مسدود شده‌اند. همان‌گونه که در مسئله‌ی~\ref{prob:chap1_causal_attribution} در فصل نخست و معادله‌ی~\ref{eq:chap4_constrained_opt} در فصل چهارم صورت‌بندی شد، چالش محوری در این بستر، یافتن بردار ماسک دودویی بهینه با نام $\mathbf{m}^* \in {\{0, 1\}}^N$ است که با اعمال بر روی $N$ جمله‌ی پرامپت ورودی ($\mathbf{x} = (x_1, \dots, x_N) \in \mathcal{S}^N$)، زیرمجموعه‌ی پیش‌آیندهای منتخب ($\mathbf{x}_S = \mathbf{x} \odot \mathbf{m}^*$) را برای تبیین زیرمجموعه‌ی هدف از جملات خروجی ($\mathbf{y}_T = {(y_j)}_{j \in T} \subseteq \mathbf{y}$) استخراج نماید:

\begin{equation}\label{eq:chap6_constrained_opt}
    \mathbf{m}^* = \arg\min_{\mathbf{m} \in {\{0, 1\}}^N} \|\mathbf{m}\|_0 \quad \text{به‌شرط آن‌که} \quad \mathcal{D}\big(f_\theta(\mathbf{x} \odot \mathbf{m}), \mathbf{y}_T\big) \le \epsilon
\end{equation}
که در آن $\|\cdot\|_0$ بیانگر اصل ضرورت و فشردگی، $\mathcal{D}$ سنجه‌ی واگرایی معنایی متناظر با اصل کفایت، و $\epsilon > 0$ آستانه‌ی خطای پذیرفتنی است. برای حل این مسئله‌ی بهینه‌سازی ترکیبیاتی بدون گرفتاری در پیچیدگی نمایی فضای حالات و بدون نیاز به استعلام‌های مکرر از مدل هدف، چارچوب جانشین مبتنی بر انرژی برای تفسیر در سطح مفهوم (\lr{ESCI}) طراحی، پیاده‌سازی و اعتبارسنجی شد. دستاوردهای این رویکرد را می‌توان در محورهای کلیدی زیر خلاصه نمود:

\paragraph{ارتقای واحد تحلیل به سطح گزاره و مهار خطای توزیعی.}
یکی از موانع ساختاری در روش‌های سنتی انتساب ویژگی در پردازش زبان طبیعی، تقلیل‌گرایی به سطح توکن‌های خرد است. دست‌کاری یا حذف توکن‌های منفرد در مدل‌های مولد خودهمبسته، گرامر متن را متلاشی کرده و منجر به بروز خطای داده‌های خارج از توزیع (\lr{OOD}) می‌شود؛ وضعیتی که در آن نوسان رفتار مدل صرفاً بازتابی از واکنش به تخریب ساختار نحوی است، نه نشانه‌ای از وابستگی محتوایی واقعی~\cite{madsen2022evaluating}. علاوه بر این، به دلیل پدیده‌ی افزونگی واژگانی، افت احتمال یک توکن خاص لزوماً به معنای محو بار معنایی آن مفهوم از خروجی نیست. در این پژوهش، با اتکا به مبانی فلسفی اصل بافتار فرگه~\cite{frege1991foundations} و نظریه‌ی شناختی ساخت‌وساز-ادغام کینتش~\cite{kintsch1998comprehension}، واحد تحلیل به سطح جمله به‌عنوان کوچک‌ترین حامل یک گزاره‌ی منطقی مستقل ارتقا یافت. این تغییر مقیاس بنیادین تضمین کرد که حذف یا گزینش مؤلفه‌های ورودی، هویت ساختاری و معنایی متن را حفظ کرده و پیش‌آیندهای استخراج‌شده به‌شکلی شفاف و انسان‌خوانا تبیین شوند.

\paragraph{مدل‌سازی چشم‌انداز انرژی جانشین بدون نیاز به نرمال‌سازی احتمالات.}
در شرایط انسداد رابط‌های تجاری، عدم‌دسترسی به بردار لوجیت‌ها مانع از به‌کارگیری معیارهای مشتق‌پذیر سنتی می‌شود، و از سوی دیگر، محاسبه‌ی تابع افراز بر روی فضای نامتناهی دنباله‌های متنی از نظر ریاضی غیرقابل‌حل است~\cite{lecun2006tutorial}. چارچوب \lr{ESCI} برای حل این معضل، یک مدل مبتنی بر انرژی (\lr{EBM}) با معماری ترنسفورمر دوسویه ($\mathcal{E}_{\text{\lr{LM}}}(\mathbf{x}, \mathbf{y}; \theta)$) را به‌عنوان ارزیاب انطباق‌پذیریِ نرمال‌نشده به‌کار می‌گیرد~\cite{bakhtin2019realfake, deng2020residual}. این شبکه با پرهیز از تولید متن (که خود منشأ بروز توهم تفسیرپذیری در مدل‌های خودتوضیح است)، تمام ظرفیت پارامتری خود را به ارزیابی یکپارچه‌ی ساختار پرامپت و پاسخ اختصاص می‌دهد. آموزش این مدل با بهره‌گیری از یک تابع زیان دوگانه مبتنی بر \lr{InfoNCE}~\cite{oord2018representation} شامل زیان وفاداری سراسری ($\mathcal{L}_{\text{\lr{fidelity}}}$) و زیان وابستگی محلی پادواقعی ($\mathcal{L}_{\text{\lr{dep}}}$) انجام می‌پذیرد. بدین‌ترتیب، مدل انرژی انطباق‌پذیری ساختاری و معنایی میان گزاره‌های ورودی و خروجی را بر مبنای متون مشاهده‌پذیر مدل هدف مدل‌سازی می‌کند، بدون آن‌که فرضی پیرامون یادگیری خمینه‌ها یا پویایی بازنمایی‌های درونی مدل جعبه‌سیاه اعمال شود.

\paragraph{استنتاج مستهلک‌شده و استخراج پیش‌آیندها در زمان‌ثابت $\mathcal{O}(1)$.}
روش‌های اختلالی کلاسیک نظیر \lr{LIME}~\cite{ribeiro2016lime} برای ارزیابی هر ورودی جدید، نیازمند کاوش در فضای ترکیبیاتی و ارسال صدها استعلام به رابط برنامه‌نویسی (\lr{API}) مدل هدف هستند که هزینه‌ی مالی و زمانی سنگینی را تحمیل می‌کند. در چارچوب \lr{ESCI}، با بهره‌گیری از پارادایم استنتاج مستهلک‌شده~\cite{tu2018learning, chen2018learning}، دانش توزیع‌شده در چشم‌انداز انرژی به یک شبکه‌ی مفسر مستقل ($\mathcal{IN}(\mathbf{x}, \mathbf{y}, T; \alpha)$) تقطیر می‌شود. مفسر با شرطی‌سازی بازنمایی‌ها بر روی شاخص‌های هدف خروجی ($T$) و مدل‌سازی انتخاب جملات به‌صورت توزیع‌های برنولی مستقل با استفاده از تابع سیگموئید، آموزش می‌بیند تا شکاف انرژی زیر را هم‌زمان با اعمال جریمه‌ی خلوتی نُرم $L_1$ بیشینه سازد:
\begin{equation}\label{eq:chap6_energy_gap}
    \Delta\mathcal{E}(\hat{\mathbf{m}}) = \mathcal{E}_{\text{\lr{LM}}}\big(\mathbf{x} \odot (\mathbf{1} - \hat{\mathbf{m}}), \mathbf{y}_T; \theta\big) - \mathcal{E}_{\text{\lr{LM}}}\big(\mathbf{x} \odot \hat{\mathbf{m}}, \mathbf{y}_T; \theta\big)
\end{equation}
این طراحی، بار محاسباتی را به فاز پیش‌آموزش منتقل کرده و در زمان استنتاج، پیش‌آیندهای اثرگذار را تنها با یک بار عبور پیش‌رو در گراف محاسباتی مفسر و در زمان‌ثابت $\mathcal{O}(1)$ بدون حتی یک استعلام از مدل هدف استخراج می‌نماید. تحلیل‌های تجربی نشان داد که نقطه‌ی سربه‌سر محاسباتی این رویکرد در مقایسه با روش‌های اختلالی پس از ممیزی تنها ۲۰ پرامپت با طول متوسط حاصل می‌شود و پس از آن، هزینه‌ی حاشیه‌ای ممیزی در \lr{ESCI} به صفر می‌رسد.

\paragraph{ارزیابی وفاداری تجربی در تقابل با خطوط پایه‌ی جعبه‌سفید و جعبه‌سیاه.}
برای سنجش جامع کیفیت انتساب، چارچوب پیشنهادی در دو مجموعه‌داده‌ی مکمل شامل \lr{HC3}~\cite{guo2023closechatgpthumanexperts} (پرسش و پاسخ آزاد مکالمه‌ای) و \lr{WikiBio}~\cite{lebret2016neural} (تولید متن مقید بر پایه‌ی داده‌های ساختاریافته) و بر روی سه مدل زبانی با معماری‌های ناهمگن شامل \lr{Llama–3.1–8B–Instruct}، \lr{Gemma–2–9B–It} و مدل تجاری مسدود \lr{GPT–4o–Mini} پیاده‌سازی و ارزیابی شد. ارزیابی وفاداری بر پایه‌ی یک خط‌لوله‌ی مداخله‌ای سه‌گانه شامل کفایت ($\mathcal{S}$)، جامعیت ($\mathcal{C}$) و شکاف وفاداری ($\Delta\mathcal{F} = \mathcal{S} - \mathcal{C}$) تحت سه عملگر محاسباتی متمایز انجام گرفت: استلزام منطقی صوری ($\mathcal{F}_{\text{\lr{NLI}}}$)، شباهت کسینوسی بیشینه‌ی جملات ($\mathcal{F}_{\text{\lr{cosine}}}$)، و لگاریتم احتمالات نرمال‌شده‌ی توکن‌ها ($\mathcal{F}_{\text{\lr{logprob}}}$)~\cite{carton2020evaluating, atanasova2023faithfulness}.

در این ارزیابی‌ها، عملکرد \lr{ESCI} هم در برابر طیف وسیعی از روش‌های جعبه‌سیاه (شامل هیئت اوراکل‌های زبانی، خودتوضیحی و درون‌نگری مدل‌های هدف، واریانت جعبه‌سیاه \lr{LIME} در سطح جمله، تطبیق رشته‌ای و خط‌پایه‌ی تصادفی خلوت دیریکله) و هم در تقابل با روش‌های دارای دسترسی ممتاز جعبه‌خاکستری و جعبه‌سفید (شامل توجه خام ترنسفورمر~\cite{vaswani2017attention}، بسط توجه~\cite{abnar2020quantifying}، گرادیان شاپ~\cite{lundberg2017unified}، ورودی ضرب‌در گرادیان~\cite{li2016visualizing, simonyan2014deepinside}، روش حذفی در سطح جمله یا \lr{Sentence-LOO}~\cite{li2017understanding, zeiler2014visualizing}، و واریانت جعبه‌خاکستری \lr{LIME}~\cite{ribeiro2016lime}) مورد سنجش قرار گرفت. 

نتایج تجربی نشان داد که اگرچه روش‌های جعبه‌سفید به دلیل دسترسی مستقیم به وزن‌ها، فعال‌سازی‌های پنهان و گرادیان‌ها از مزیتی ساختاری در سنجه‌های مبتنی بر افت احتمالات توکن‌ها برخوردارند، اما در سنجش وفاداری معنایی بر مبنای استلزام منطقی صوری ($\Delta\mathcal{F}_{\text{\lr{NLI}}}$)، مفسر جعبه‌سیاه \lr{ESCI} با ثبت شکاف‌های مثبت $0.108 \pm 0.004$ بر روی \lr{Llama–3.1} و $0.228 \pm 0.004$ بر روی \lr{Gemma–2}، عملکردی فراتر از مکانیزم‌های توجه درونی (شکاف $0.093$ و $0.084$) و ورودی ضرب‌در گرادیان (شکاف $-0.027$ و $+0.029$) به نمایش گذاشت و به روش‌های پرهزینه‌ی اختلالی نظیر \lr{Sentence-LOO} (شکاف $0.123$ و $0.295$) نزدیک شد. افزون بر این، در ارزیابی معقولیت انتساب بر روی \lr{WikiBio}، مفسر پیشنهادی هم‌راستایی رتبه‌ای بالایی با روش‌های جریان توجه (نمره‌ی \lr{nDCG} معادل $0.726$ با بسط توجه و $0.743$ با توجه خام) ثبت کرد که نشان می‌دهد چشم‌انداز انرژی سرشکن‌شده قادر است ساختار کلان وابستگی‌های گزاره‌ای را بدون نیاز به بازرسی لایه‌های درونی بازتاب دهد.

\paragraph{مرزبندی معرفت‌شناختی و سرنخ‌های عینی اثرگذاری.}
در نهایت، تأکید بر این نکته ضروری است که انتساب‌های استخراج‌شده در چارچوب \lr{ESCI} بر پایه‌ی الگوهای انطباق‌پذیری متنی در متون مشاهده‌پذیر شکل گرفته‌اند. به‌بیانی دیگر، پیش‌آیندهای شناسایی‌شده شواهد و نشانگرهایی تجربی و عینی از نحوه‌ی تأثیر اجزای پرامپت بر گزاره‌های خروجی فراهم می‌آورند که از نظر تجربی با خطوط زیربنایی تأثیرگذاری در مدل همسو هستند. نگریستن به مسئله‌ی انتساب جعبه‌سیاه از این زاویه، بستری عملی و محاسبه‌پذیر برای پژوهش‌های آتی در حوزه‌ی ممیزی مدل‌های مولد بنا می‌نهد.


\section{محدودیت‌های پژوهش}\label{sec:6-2}

علی‌رغم دستاوردهای روش‌شناختی و مزایای محاسباتی چارچوب \lr{ESCI}، استقرار و کاربست این متدولوژی با محدودیت‌های ساختاری و واقعیت‌های عینی حاکم بر محیط‌های جعبه‌سیاه روبه‌رو است که تبیین شفاف آن‌ها برای هدایت پژوهش‌های آینده الزامی است:

\paragraph{فقدان ذاتی حقیقت پایه و قلمرو رفتاری.}
ارزیابی وفاداری توضیحات در پردازش زبان طبیعی و مدل‌های مولد، به‌طور بنیادین با چالش فقدان حقیقت پایه‌ی قطعی و جهان‌شمول مواجه است~\cite{schneider2024explainable, moran2026towardinterpretable}. از آنجا که مسیرهای محاسباتی خودهمبسته در مدل‌های مسدود تجاری غیرقابل‌مشاهده و مات هستند، نه حاشیه‌نویسی‌های انسانی، نه نگاشت‌های مکاشفه‌ای مبتنی بر تطبیق واژگانی (نظیر معیارهای اعتبارسنجی خارجی \lr{AIS}~\cite{rashkin2023measuring} یا خط‌پایه‌ی تطبیق رشته‌ای)، و نه حتی سنجه‌های مداخله‌ای پادواقعی، هیچ‌یک استانداردی مطلق برای انتساب ارائه نمی‌دهند. این چالش، بازتاب‌دهنده‌ی یک محدودیت ساختاری در ماهیت ممیزی جعبه‌سیاه در پردازش زبان طبیعی است، و یک نقص الگوریتمی گذرا که با راهکارهای ساده برطرف شود نیست.

از یک سو، همان‌گونه که در ادبیات سنجه‌های تفسیرپذیری تبیین شده است، معیارهای مداخله‌ای می‌توانند تحت تأثیر قانون گودهارت دچار سوگیری ارزیابی شوند~\cite{ju2022logictraps, hsia2024goodhart}؛ وضعیتی که در آن مفسرها به‌جای انتساب گزاره‌های دارای نقش ساختاری، به‌سوی انتخاب بخش‌هایی سوق می‌یابند که بیشترین شوک توزیعی را در فضای خروجی ایجاد کنند. از سوی دیگر، یافته‌های تجربی این پژوهش در هم‌سویی با مطالعات کاماهی و یعقوب‌زاده~\cite{kamahi2024counterfactuals} نشان داد که مدل‌های زبانی در هنگام حذف جملات ورودی در آزمایش‌های جامعیت، به‌طور گسترده به دانش پارامتریک درونی خود تکیه می‌کنند؛ پدیده‌ای که مانع از سقوط کامل نمره‌ی جامعیت ($\mathcal{C}_{\text{\lr{omp}}}$) به صفر شده و مقادیر آن را در دامنه‌ی $0.42$ تا $0.60$ تثبیت می‌کند. در نتیجه‌ی این واقعیت‌های ساختاری، چارچوب \lr{ESCI} یک سازوکار انتساب رفتاری بر پایه‌ی الگوهای همبستگی ساختاری ورودی-خروجی ارائه می‌دهد و نباید آن را با بازرسی مدارهای مکانیکی در لایه‌های پنهان شبکه هم‌ارز دانست~\cite{elhage2021mathematical, friedman2023comparing}. بر این اساس، به‌جای جستجوی یک حقیقت پایه‌ی دست‌نیافتنی، گسترش ارزیابی‌ها در بسترهای عملیاتی متنوع‌تر برای سنجش پایداری انتساب تحت شرایط گوناگون دنیای واقعی ضروری خواهد بود.

\paragraph{سربار محاسباتی پیش‌آموزش در برابر استنتاج مستهلک‌شده.}
اگرچه چارچوب \lr{ESCI} با دستیابی به پیچیدگی زمان‌ثابت $\mathcal{O}(1)$ نیاز به استعلام‌های مکرر در زمان استنتاج را به صفر می‌رساند، اما این کارایی عملیاتی از طریق انتقال بار پردازشی به فاز پیش‌آموزش حاصل می‌شود. تولید دادگان جفت‌های پرامپت-پاسخ ($20,000$ نمونه و حتی بیشتر)، ساخت نمونه‌های منفی تقابلی، و اجرای فرایند لنگراندازی دوره‌ای برای همگام نگه‌داشتن مدل انرژی با مفسر، مستلزم صرف هزینه و زمان اولیه در تعامل با \lr{API} مدل هدف است. اگرچه این هزینه پس از ممیزی تعداد محدودی پرامپت (نقطه‌ی سربه‌سر در مقایسه با \lr{LIME}) کاملاً مستهلک می‌شود، اما برای کاربردهایی که با وظایف گذرا و تک‌نمونه‌ای سروکار دارند یا کاربرانی با محدودیت‌های پردازشی شدید، پرداخت این هزینه‌ی اولیه‌ی آماده‌سازی نیازمند برنامه‌ریزی منابع است.

\paragraph{تخصیص وظیفه‌محور و محدودیت ظرفیت مدل‌های جانشین فشرده.}
به‌منظور امکان‌پذیر ساختن آموزش و استقرار چارچوب بر روی سخت‌افزارهای متعارف و در دسترس (نظیر دو پردازنده‌ی گرافیکی \lr{NVIDIA T4})، مدل‌های جانشین انرژی و مفسر در این پژوهش با ساختاری سبک و ظرفیت پارامتری فشرده (مجموعاً ۱۸۱ میلیون پارامتر که تنها حدود ۷۱ میلیون آن در مفسر آموزش‌پذیر است) پیکربندی شدند. این طراحی متمرکز، اگرچه عملکرد مناسبی در یادگیری الگوهای هم‌نشینی در دامنه‌های آزمون (مانند مکالمات روزمره‌ی \lr{HC3} و بیوگرافی‌های \lr{WikiBio}) به‌همراه دارد، اما بازنمایی‌های آموخته‌شده را وابسته به دامنه‌ی داده‌های آموزشی می‌سازد. مطابق با مشاهدات بختیین و همکاران~\cite{bakhtin2019realfake} پیرامون افت دقت مدل‌های انرژی در تعمیم بین‌دامنه‌ای در غیاب داده‌های ترکیبی، مفسری که منحصراً بر روی متون عمومی آموزش دیده است، بدون گذراندن فاز پیش‌آموزش تکمیلی یا ارتقای ظرفیت پارامتری، قابلیت تعمیم مستقیم به متون تخصصی (نظیر اسناد حقوقی پیچیده یا متون پزشکی تخصصی) را نخواهد داشت.

\paragraph{کران‌های تفکیک‌پذیری در سطح جمله.}
انتخاب جمله به‌عنوان مفهوم اتمیک و واحد عملیاتی تحلیل، اگرچه از فروپاشی نحوی و افتادن در دام شوک‌های توزیعی ناشی از اختلالات توکن‌محور جلوگیری می‌کند، اما به‌طور ساختاری حد بالایی برای تفکیک‌پذیری معنایی تعیین می‌نماید. در متون زبانی پیچیده، یک جمله‌ی مرکب و طولانی ممکن است شامل چندین گزاره‌ی فرعی، قیود شرطی و اطلاعات توصیفی متمایز باشد. در ساختار فعلی، مفسر کل جمله را به‌صورت یک گزاره‌ی یکپارچه ارزیابی می‌کند که این امر تفکیک دقیق‌تر سهم هر یک از عبارات درون‌جمله‌ای را ناممکن می‌سازد.


\section{پیشنهادها برای تحقیقات آتی}\label{sec:6-3}

بر پایه‌ی دستاوردهای به‌دست‌آمده، اعتبارسنجی‌های تجربی و محدودیت‌های شناسایی‌شده، مسیرهای پژوهشی زیر برای گسترش و تکامل این متدولوژی پیشنهاد می‌گردد:

\paragraph{مقیاس‌پذیری در پیکره‌های چندوظیفه‌ای و بافتارهای طولانی.}
گام نخست در جهت ارتقای تعمیم‌پذیری چارچوب \lr{ESCI}، افزایش مقیاس پارامتری شبکه‌های جانشین و آموزش آن‌ها بر روی پیکره‌های چندوظیفه‌ای و متون چنددامنه‌ای گسترده است تا مفسر بتواند قابلیت تعمیم‌پذیری بدون نیاز به تنظیم مجدد را در مواجهه با پرامپت‌های ناآشنا به‌دست آورد. علاوه بر این، در رویارویی با اسناد بسیار طولانی در سیستم‌های تولید تقویت‌شده با بازیابی (\lr{RAG})، بازطراحی ماژول‌های توجه متقاطع با استفاده از مکانیزم‌های توجه سلسله‌مراتبی، امکان ارزیابی دوگانه در دو سطح کلان (پاراگراف‌ها و بخش‌ها) و خرد (جملات) را فراهم خواهد ساخت تا مقیاس‌پذیری محاسباتی برای پنجره‌های بافتار وسیع تضمین گردد.

\paragraph{فشرده‌سازی خودکار پرامپت.}
فراتر از کاربردهای ممیزی و تفسیرپذیری، توانایی استخراج سریع پیش‌آیندهای اثرگذار، ظرفیت بالایی برای بهینه‌سازی مستقیم سامانه‌های زبانی ارائه می‌دهد. ماسک‌های خلوت تولیدشده توسط شبکه‌ی مفسر می‌توانند به‌عنوان مبنایی برای فشرده‌سازی خودکار پرامپت پیش از ارسال به مدل‌های تجاری به‌کار روند. با حذف پیشینی جملات فاقد ارزش انتسابی از بافتار، طول پرامپت ارسالی به رابط‌های برنامه‌نویسی کاهش می‌یابد که این امر مستقیماً به کاهش تأخیر استنتاج و افت چشمگیر هزینه‌های مالی در سیستم‌های بلادرنگ منجر خواهد شد.

\paragraph{ممیزی و اعتبارسنجی استدلال در زنجیره‌ی افکار.}
در استدلال‌های گام‌به‌گام و چندمرحله‌ای مدل‌های زبانی نظیر پرامپت‌گذاری زنجیره‌ی افکار~\cite{wei2022chain}، مدل‌ها مستعد افتادن در دام استدلال هدفمند، توجیهات نادرست و توهمات متنی هستند~\cite{turpin2023language}. چشم‌انداز انرژی چارچوب \lr{ESCI} می‌تواند به‌عنوان یک ناظر و ارزیاب مستقل بر زنجیره‌ی استدلال عمل کند. در این الگو، با ارزیابی انطباق‌پذیری انرژی میان هر گام میانی با گزاره‌های ماقبل خود، مراحلی که دارای عدم‌انطباق بالایی در چشم‌انداز انرژی هستند شناسایی شده و به‌عنوان خطاهای استدلالی پیش از تولید پاسخ نهایی فیلتر یا تصحیح می‌گردند.

\paragraph{توسعه به مدل‌های بینایی-زبانی و دامنه‌ی چندوجهی.}
اصول ساختاری حاکم بر چارچوب \lr{ESCI} به زبان متنی محدود نبوده و قابل‌تعمیم به مدل‌های چندوجهی است. با جایگزینی واحد جمله با بخش‌های ساختاریافته‌ی تصویر یا اجزای قطعه‌بندی‌شده‌ی معنایی در مدل‌های بینایی-زبانی، می‌توان یک چشم‌انداز انرژی مشترک میان نواحی بصری و گزاره‌های متنی آموزش داد.

\paragraph{ارزیابی‌های تعاملی و سنجش اثرگذاری بر کاربر انسان.}
پیرو مبانی نظری دوشی-ولز و کیم~\cite{doshivelez2017rigorous} مبنی بر لزوم ارزیابی‌های کاربردمحور در تفسیرپذیری، گام ضروری بعدی، اجرای مطالعات بالینی با حضور کاربران انسانی در حلقه‌ی سیستم (حوزه‌ی تعامل انسان و رایانه) است. سنجش این امر که آیا ارائه‌ی پیش‌آیندهای گزاره‌ایِ استخراج‌شده توسط \lr{ESCI} در عمل سرعت، دقت و اعتماد کالیبره‌شده‌ی تحلیل‌گران را در فرایندهای ممیزی، تصحیح خطاها و اعتبارسنجی اسناد افزایش می‌دهد یا خیر، معیاری تعیین‌کننده برای میزان اثربخشی این متدولوژی در محیط‌های عملیاتی دنیای واقعی خواهد بود.

\vspace{0.5cm}
در پایان، این پایان‌نامه با پایه‌ریزی چارچوب \lr{ESCI}، گامی در جهت عبور از رویکردهای تقلیل‌گرای توکن‌محور به‌سوی تفسیرپذیری محاسباتی و معنادار در سطح مفاهیم برداشت. با تلفیق مدل‌سازی انرژی در سطح گزاره‌ها و تقطیر دانش در یک مفسر مستهلک‌شده، این پژوهش نشان داد که می‌توان رفتار مدل‌های زبانی عظیم و مسدود را بدون نقض محدودیت‌های دسترسی جعبه‌سیاه، با مقداری خطا در زمان‌ثابتِ استنتاج ممیزی کرد. امید است که پایه‌های نظری و تجربی استوارشده در این رساله، افق‌های نوینی را برای شفافیت ساختاری و استقرار مسئولانه‌ی نسل‌های آینده‌ی هوش مصنوعی مولد بگشاید.